\documentclass[11pt]{article}

\usepackage[a4paper,margin=1in]{geometry}
\usepackage[T1]{fontenc}
\usepackage[utf8]{inputenc}
\usepackage{lmodern}
\usepackage{microtype}
\usepackage{amsmath,amssymb,amsthm,mathtools}
\usepackage{xcolor}
\usepackage[colorlinks=true,linkcolor=blue!45!black,
  citecolor=blue!45!black,urlcolor=blue!45!black]{hyperref}

\newtheorem{theorem}{Theorem}[section]
\newtheorem{proposition}[theorem]{Proposition}
\newtheorem{corollary}[theorem]{Corollary}
\theoremstyle{definition}
\newtheorem{example}[theorem]{Example}
\theoremstyle{remark}
\newtheorem{remark}[theorem]{Remark}

\newcommand{\Pcal}{\mathcal P}
\newcommand{\Mcal}{\mathcal M}
\newcommand{\Lc}{\mathcal L_c}
\newcommand{\PiC}{\Pi}
\newcommand{\dd}{\,\mathrm d}
\newcommand{\Id}{\mathrm{Id}}

\title{Optimal Transport as the Maximal Convex Extension of a Cost}
\author{Gabriel Peyr\'e}
\date{July 2026}

\begin{document}
\maketitle

\begin{abstract}
The Kantorovich cost is often described as the convex extension of the
pointwise cost from pairs of Dirac masses to pairs of probability measures.
The precise statement contains an essential word: it is the \emph{largest}
weakly lower-semicontinuous jointly convex extension. Joint convexity,
weak-star continuity and agreement on Dirac masses do not by themselves imply
uniqueness. This note states the sharp probability-measure result, gives its
short barycentric proof, exhibits a two-point counterexample to bare
uniqueness, and records the measure-cone theorem of Savar\'e and Sodini.
\end{abstract}

\section{The probability-measure statement}

Let $\mathcal X$ and $\mathcal Y$ be compact metric spaces and let
$c\in C(\mathcal X\times\mathcal Y)$ be finite. Endow
$\Pcal(\mathcal X)$ and $\Pcal(\mathcal Y)$ with their narrow, equivalently
weak-star, topologies. For $\alpha\in\Pcal(\mathcal X)$ and
$\beta\in\Pcal(\mathcal Y)$, write $\PiC(\alpha,\beta)$ for the probability
measures on $\mathcal X\times\mathcal Y$ with marginals $\alpha$ and $\beta$,
and define
\begin{equation}\label{eq:kantorovich-cost}
  \Lc(\alpha,\beta)
  \coloneqq
  \min_{\pi\in\PiC(\alpha,\beta)}
  \int_{\mathcal X\times\mathcal Y}c(x,y)\dd\pi(x,y).
\end{equation}
Compactness and continuity ensure that the minimum is attained.

The geometry behind the result is elementary. Introduce the Dirac embedding
\[
  \jmath:\mathcal X\times\mathcal Y
  \longrightarrow \Pcal(\mathcal X)\times\Pcal(\mathcal Y),
  \qquad
  \jmath(x,y)=(\delta_x,\delta_y).
\]
A probability measure $\pi$ on $\mathcal X\times\mathcal Y$ belongs to
$\PiC(\alpha,\beta)$ exactly when the barycenter of
$\jmath_\sharp\pi$ is $(\alpha,\beta)$, in the weak sense
\begin{equation}\label{eq:barycenter-dirac-pairs}
  (\alpha,\beta)
  =\int_{\mathcal X\times\mathcal Y}
  (\delta_x,\delta_y)\dd\pi(x,y).
\end{equation}
Thus couplings are precisely the convex decompositions of the pair of
marginals into compatible pairs of extreme points.

\begin{theorem}[Maximal convex-extension characterization]
\label{thm:maximal-convex-extension}
The functional $\Lc$ has the following properties.
\begin{enumerate}
  \item It is jointly convex and weak-star continuous on
  $\Pcal(\mathcal X)\times\Pcal(\mathcal Y)$.
  \item It extends the pointwise cost:
  \[
    \Lc(\delta_x,\delta_y)=c(x,y)
    \qquad\text{for every }(x,y)\in\mathcal X\times\mathcal Y.
  \]
  \item If $F:\Pcal(\mathcal X)\times\Pcal(\mathcal Y)
  \to(-\infty,+\infty]$ is proper, jointly convex and weak-star lower
  semicontinuous, and
  \begin{equation}\label{eq:dirac-domination}
    F(\delta_x,\delta_y)\leq c(x,y)
    \qquad\text{for every }(x,y),
  \end{equation}
  then
  \begin{equation}\label{eq:maximality}
    F(\alpha,\beta)\leq\Lc(\alpha,\beta)
    \qquad\text{for every }(\alpha,\beta).
  \end{equation}
\end{enumerate}
Consequently, $\Lc$ is the pointwise largest jointly convex weak-star
lower-semicontinuous functional dominated by $c$ on Dirac pairs. Since $\Lc$
is continuous and has equality on Dirac pairs, it is also the largest member
of the smaller class specified by the three properties in the question.
\end{theorem}

\begin{proof}
The only coupling of $\delta_x$ and $\delta_y$ is
$\delta_{(x,y)}$, which proves the Dirac identity. If
$\pi_i\in\PiC(\alpha_i,\beta_i)$ are optimal and $t\in[0,1]$, then
$t\pi_0+(1-t)\pi_1$ couples
$t\alpha_0+(1-t)\alpha_1$ and $t\beta_0+(1-t)\beta_1$.
Its cost is the corresponding convex combination, proving joint convexity.

We briefly verify continuity. Suppose
$(\alpha_n,\beta_n)\to(\alpha,\beta)$. Any convergent subsequence of optimal
plans $\pi_n$ has a limit in $\PiC(\alpha,\beta)$; continuity of $c$ then gives
\[
  \Lc(\alpha,\beta)
  \leq \liminf_n \Lc(\alpha_n,\beta_n).
\]
Conversely, start from an optimal $\pi\in\PiC(\alpha,\beta)$. On compact metric
spaces, weak convergence is equivalent to convergence in the Wasserstein
distance associated with any bounded compatible metric. Couple $\alpha_n$ to
$\alpha$ and $\beta_n$ to $\beta$ with costs tending to zero, and glue these
two couplings through $\pi$. This produces
$\widetilde\pi_n\in\PiC(\alpha_n,\beta_n)$ with
$\widetilde\pi_n\rightharpoonup\pi$. Uniform continuity of $c$ gives
\[
  \limsup_n\Lc(\alpha_n,\beta_n)
  \leq\lim_n\int c\dd\widetilde\pi_n
  =\int c\dd\pi
  =\Lc(\alpha,\beta).
\]

It remains to prove maximality. Fix any $\pi\in\PiC(\alpha,\beta)$.
Equations~\eqref{eq:barycenter-dirac-pairs} and~\eqref{eq:dirac-domination},
together with Jensen's inequality, yield
\[
  F(\alpha,\beta)
  \leq
  \int F(\delta_x,\delta_y)\dd\pi(x,y)
  \leq
  \int c(x,y)\dd\pi(x,y).
\]
For completeness, this version of Jensen follows by approximating $\pi$ by
finitely supported probability measures, applying finite-dimensional
convexity, and passing to the limit using lower semicontinuity. Taking the
infimum over $\pi$ proves~\eqref{eq:maximality}.
\end{proof}

Equivalently, the theorem gives the convex-roof formula
\begin{equation}\label{eq:convex-roof}
  \Lc(\alpha,\beta)
  =
  \inf\left\{
    \int c(x,y)\dd\pi(x,y):
    (\alpha,\beta)=\int(\delta_x,\delta_y)\dd\pi(x,y)
  \right\}.
\end{equation}
The marginal constraints in optimal transport are therefore exactly the
barycenter constraint in the convexification of the Dirac data.

\section{Why the three properties do not imply uniqueness}

The maximality clause cannot be omitted, even when all spaces are finite and
``weak-star continuous'' means ordinary continuity.

\begin{example}[A two-point counterexample]
Let $\mathcal X=\mathcal Y=\{0,1\}$ and $c(x,y)=|x-y|$. Write
\[
  \alpha_p=(1-p)\delta_0+p\delta_1,
  \qquad
  \beta_q=(1-q)\delta_0+q\delta_1.
\]
Then
\[
  \Lc(\alpha_p,\beta_q)=|p-q|.
\]
The alternative functional
\[
  F(\alpha_p,\beta_q)=(p-q)^2
\]
is continuous and jointly convex because it is the square of a linear
functional. For the four Dirac pairs $p,q\in\{0,1\}$, it satisfies
$F(\delta_p,\delta_q)=|p-q|=c(p,q)$. Nevertheless,
\[
  F(\alpha_{1/2},\beta_0)=\frac14
  <\frac12=\Lc(\alpha_{1/2},\beta_0).
\]
Hence joint convexity, weak-star continuity and exact Dirac values imply only
$F\leq\Lc$, not $F=\Lc$.
\end{example}

One may turn the maximality statement into a literal uniqueness axiom in
either of two equivalent ways: require $F$ to be pointwise maximal among the
admissible convex extensions, or require $F$ to equal the convex roof of its
Dirac values. Either condition forces $F=\Lc$ by
Theorem~\ref{thm:maximal-convex-extension}.

\section{The Savar\'e--Sodini measure-cone theorem}

The exact theorem of Savar\'e and Sodini is more general and cleaner from the
viewpoint of convex analysis~\cite[Theorem~4.4]{SavareSodini2022}. Let
$\Mcal(\mathcal X)$ denote the vector space of finite signed Radon measures,
with weak topology induced by $C_b(\mathcal X)$. Define the singular
functional on $\Mcal(\mathcal X)\times\Mcal(\mathcal Y)$ by
\begin{equation}\label{eq:singular-dirac-functional}
  \mathsf F_c(\mu,\nu)
  \coloneqq
  \begin{cases}
    m c(x,y),
      & (\mu,\nu)=(m\delta_x,m\delta_y),\quad m\geq0,\\
    +\infty, & \text{otherwise}.
  \end{cases}
\end{equation}
with the convention $0\cdot(+\infty)=0$. In the formula below,
$\PiC(\mu,\nu)$ denotes the nonnegative finite measures on
$\mathcal X\times\mathcal Y$ with marginals $\mu$ and $\nu$. Extend the
transport functional to the same vector space by
\begin{equation}\label{eq:cone-transport-functional}
  \mathsf T_c(\mu,\nu)
  \coloneqq
  \begin{cases}
    \displaystyle\min_{\gamma\in\PiC(\mu,\nu)}\int c\dd\gamma,
      & \mu,\nu\geq0,\quad \mu(\mathcal X)=\nu(\mathcal Y),\\
    +\infty, & \text{otherwise}.
  \end{cases}
\end{equation}

\begin{theorem}[Savar\'e--Sodini]
Let $\mathcal X$ and $\mathcal Y$ be completely regular spaces and let
$c:\mathcal X\times\mathcal Y\to[0,+\infty]$ be lower semicontinuous and
proper. Then $\mathsf T_c$ is the weakly lower-semicontinuous convex envelope
of $\mathsf F_c$. Equivalently, it is the largest weakly lower-semicontinuous
convex functional $G$ such that $G\leq\mathsf F_c$. It is, in addition,
sublinear and satisfies
\[
  \mathsf T_c(m\delta_x,m\delta_y)=m c(x,y).
\]
\end{theorem}

On compact metric spaces, the proof is the same barycentric argument as
above. If $\mu$ and $\nu$ have common mass $m>0$ and
$\gamma\in\PiC(\mu,\nu)$, then
\[
  (\mu,\nu)
  =\int_{\mathcal X\times\mathcal Y}
  (m\delta_x,m\delta_y)\dd\frac{\gamma}{m}(x,y).
\]
Thus every convex $G\leq\mathsf F_c$ satisfies
\[
  G(\mu,\nu)
  \leq\int G(m\delta_x,m\delta_y)\dd\frac{\gamma}{m}
  \leq\int c\dd\gamma,
\]
and hence $G\leq\mathsf T_c$. Convexity and lower semicontinuity of
$\mathsf T_c$ give the reverse envelope inclusion. The general completely
regular case requires the compactness theorem for plans with converging
marginals established in~\cite{SavareSodini2022}.

There is also a one-line dual interpretation. The convex conjugate of
$\mathsf F_c$ is the indicator of
\[
  \left\{(f,g)\in C_b(\mathcal X)\times C_b(\mathcal Y):
  f(x)+g(y)\leq c(x,y)\right\}.
\]
Fenchel--Moreau therefore identifies the biconjugate of $\mathsf F_c$ with
the Kantorovich dual, while the envelope theorem identifies the same
biconjugate with $\mathsf T_c$. This is the convex-analytic core of
Kantorovich duality in the Savar\'e--Sodini proof.

\begin{thebibliography}{1}

\bibitem{SavareSodini2022}
G.~Savar\'e and G.~E. Sodini.
\newblock A simple relaxation approach to duality for optimal transport
  problems in completely regular spaces.
\newblock \emph{Journal of Convex Analysis}, 29(1):1--12, 2022.
\newblock \url{https://cvgmt.sns.it/paper/5594/}.

\end{thebibliography}

\end{document}
